\documentclass[a4paper,10pt]{article}

% Set \withphototrue and provide me.jpg to enable the optional photo.
\newif\ifwithphoto
\withphotofalse

\usepackage[utf8]{inputenc}
\usepackage[T1]{fontenc}
\usepackage[left=0.65in,right=0.65in,top=0.55in,bottom=0.55in]{geometry}
\usepackage{lmodern}
\usepackage{microtype}
\usepackage[hidelinks]{hyperref}
\usepackage{enumitem}
\usepackage{titlesec}
\usepackage{graphicx}
\input{glyphtounicode}
\pdfgentounicode=1
\hypersetup{pdftitle={Peng Li -- CV -- Acoustic Modelling, CFD and FSI},pdfauthor={Peng Li}}
\pagestyle{empty}
\setlength{\parindent}{0pt}
\setlength{\parskip}{0pt}
\urlstyle{same}
\titleformat{\section}{\normalsize\bfseries}{}{0pt}{}[\vspace{4pt}\titlerule]
\titlespacing*{\section}{0pt}{7pt}{4pt}
\setlist[itemize]{leftmargin=1.25em,itemsep=2pt,topsep=2pt,parsep=0pt,partopsep=0pt}
\setlist[enumerate]{leftmargin=1.5em,itemsep=4pt,topsep=2pt,parsep=0pt}
\setlength{\emergencystretch}{2em}
\newcommand{\contact}{Stockholm, Sweden \quad | \quad +46 79 010 31 96 \quad | \quad
\href{mailto:pengli.2810@gmail.com}{pengli.2810@gmail.com}}
\newcommand{\specialism}{Acoustic Modelling | Computational Fluid Dynamics | Fluid--Structure Interaction}

\begin{document}
\ifwithphoto
\noindent
\begin{minipage}[c]{0.78\textwidth}
{\LARGE\textbf{Peng Li}}\\[4pt]
\textbf{Acoustic Modelling | CFD | Fluid--Structure Interaction}\\[5pt]
\contact
\end{minipage}\hfill
\begin{minipage}[c]{0.18\textwidth}
\raggedleft\includegraphics[width=2.65cm,height=3.05cm,keepaspectratio]{me.jpg}
\end{minipage}
\else
\begin{center}
{\LARGE\textbf{Peng Li}}\\[4pt]
\textbf{\specialism}\\[4pt]
\contact\\[3pt]
\href{https://orcid.org/0009-0006-0575-0865}{ORCID: 0009-0006-0575-0865}
\end{center}
\vspace{-4pt}
\fi

\section{Professional Summary}
Ph.D. researcher at KTH Royal Institute of Technology specialising in computational acoustics,
wave propagation, and coupled fluid--structure interaction (FSI). Develop analytical acoustic
models and high-fidelity CFD--FEM simulations to explain sound generation, duct resonances,
and flow-induced vibration. Experience includes compressible large eddy simulation (LES),
acoustic-source decomposition, spectral analysis, scientific programming, and simulations on
Swedish HPC infrastructure. Bring a foundation in applied mathematics, fluid mechanics, and
structural dynamics relevant to acoustic modelling and signature analysis in underwater applications.
\textbf{Ph.D. defence scheduled for 4 December 2026.}

\section{Technical Skills}
\textbf{Acoustic modelling:} Convected wave equations; duct acoustics; acoustic resonances;
mode matching; evanescent modes; junction impedance; acoustic--structural coupling;
Ffowcs Williams--Hawkings (FW--H) analogy; duct Green's functions.\\[3pt]
\textbf{CFD and FSI:} Compressible LES; finite-volume and finite-element methods;
strong two-way coupling; overset and deforming meshes; hyperelasticity;
flow-induced vibration; transient compressible flow and shock-wave propagation.\\[3pt]
\textbf{Signal and modal analysis:} FFT; power spectral density; cross-spectra;
coherence and phase; spectral proper orthogonal decomposition (SPOD);
wavenumber--frequency analysis; structural eigenmodes.\\[3pt]
\textbf{Scientific computing:} Python, MATLAB, C/C++, Fortran, Mathematica;
parameter studies; model verification and validation; mesh and time-step convergence;
HPC simulation workflows (NAISS resources at PDC and NSC).\\[3pt]
\textbf{Engineering software:} Simcenter STAR-CCM+, OpenFOAM, COMSOL Multiphysics,
ABAQUS, ANSYS, Tecplot, Gnuplot, Inkscape, LaTeX.

\section{Research and Professional Experience}
\textbf{Ph.D. Researcher -- Computational Aeroacoustics and FSI}
\hfill Mar 2022 -- Present\\
\textit{Department of Engineering Mechanics, FLOW, KTH Royal Institute of Technology, Stockholm}
\begin{itemize}
\item Develop three-dimensional upper-airway models coupling compressible LES with
FEM-based soft-tissue mechanics in STAR-CCM+, using strong two-way FSI,
Neo-Hookean materials, and moving overset meshes.
\item Develop an analytical wave-propagation and resonance model for a bifurcated duct
with a compliant splitter, combining convected plane waves, mode-matched junction
impedance, and frequency-dependent structural compliance.
\item Verify acoustic-model predictions against COMSOL eigenfrequency calculations:
frequency differences within 0.002\% for 27 rigid-splitter resonances and 1.03\% for
the first ten soft-splitter common resonances in the tested geometry
(unpublished manuscript).
\item Evaluate monopole, dipole, and quadrupole sound-source contributions using the
FW--H analogy and duct Green's functions; relate acoustic pressure spectra to
unsteady loading, structural motion, and breathing conditions.
\item Use pressure/vorticity SPOD, wavenumber--frequency analysis, cross-spectra,
coherence, and modal analysis to distinguish acoustic standing waves from convective
flow structures and identify fluid--structure coupling mechanisms.
\item Conduct mesh and time-step convergence studies and simulations using NAISS
HPC resources at PDC and NSC; lead analysis, visualisation, and manuscript writing
in collaboration with researchers in fluid mechanics, FSI, and acoustics.
\end{itemize}

\vspace{3pt}
\textbf{Research Assistant -- Transient Modelling and Compressible CFD}
\hfill Dec 2018 -- Mar 2020\\
\textit{Joyson Safety Systems, Aschaffenburg, Germany}
\begin{itemize}
\item Developed Python models of pyrotechnical initiators from mass and energy
conservation, combustion, heat transfer, and nozzle-flow relations.
\item Coupled 0-D model outputs to 3-D OpenFOAM simulations to resolve transient
shock-wave propagation and assess airbag-inflator performance and structural integrity.
\item Compared predicted pressure histories with closed-bomb measurements and
performed parameter-sensitivity studies to support product development;
work published in \textit{Journal of Energy and Power Technology}.
\end{itemize}

\newpage

\section{Education}
\textbf{Ph.D. studies -- Computational Fluid Dynamics}
\hfill Mar 2022 -- Present\\
KTH Royal Institute of Technology, Stockholm, Sweden\\
Thesis: \textit{Computational Aero-acoustics of Snoring with Impact on Obstructive Sleep Apnea Diagnosis}\\
\textbf{Defence scheduled: 4 December 2026.}\\
Supervision: Mihai Mihaescu, Marco Laudato, and Susann Boij.

\vspace{5pt}
\textbf{M.Sc. in Mechanics (Mechanik)} \hfill Oct 2016 -- Jan 2020\\
Technical University of Darmstadt, Germany\\
Final grade: \textbf{1.33 (very good)}; thesis: \textbf{1.0} (German grading system).\\
Thesis: \textit{Analysis and modeling of ballistic behaviour of pyrotechnical initiators
and their impact in airbag inflators applications}.

\vspace{5pt}
\textbf{B.Eng. in Engineering Mechanics} \hfill Sep 2011 -- Jun 2015\\
Dalian University of Technology, China\\
Overall mark: 86/100; thesis mark: 98/100.

\section{Selected Peer-Reviewed Journal Publications}
\begin{enumerate}
\item \textbf{Li, P.}, Laudato, M., and Mihaescu, M. (2026).
Computational analysis of palatal non-apneic snoring sound generation using a simplified
human upper-airway model. \textit{Physics of Fluids}, \textbf{38}(8), 081909.
\textbf{Featured Article}.
\href{https://doi.org/10.1063/5.0346307}{doi:10.1063/5.0346307}.
\item \textbf{Li, P.}, Laudato, M., and Mihaescu, M. (2023).
Time-dependent fluid-structure interaction simulations of a simplified human soft palate.
\textit{Bioengineering}, \textbf{10}(11), 1313.
\href{https://doi.org/10.3390/bioengineering10111313}{doi:10.3390/bioengineering10111313}.
\item \textbf{Li, P.}, Chang, C.-Y., and Sadiki, A. (2021).
Numerical investigation of pyrotechnical initiators and their impact in airbag inflator
applications: Ballistic modeling and functional analysis.
\textit{Journal of Energy and Power Technology}, \textbf{3}(2), article 026.
\href{https://doi.org/10.21926/jept.2102026}{doi:10.21926/jept.2102026}.
\item Wang, B., \textbf{Li, P.}, and Li, R. (2016).
Symplectic superposition method for new analytic buckling solutions of rectangular thin plates.
\textit{International Journal of Mechanical Sciences}, \textbf{119}, 432--441.
\href{https://doi.org/10.1016/j.ijmecsci.2016.11.006}{doi:10.1016/j.ijmecsci.2016.11.006}.
\end{enumerate}

\section{Unpublished Manuscripts}
\begin{enumerate}
\item \textbf{Li, P.}, Laudato, M., Boij, S., and Mihaescu, M.
\textit{Acoustic resonances of a bifurcated duct with a compliant splitter}.
Analytical wave modelling, junction impedance, acoustic--structural coupling, and FEM verification.
\item \textbf{Li, P.}, Laudato, M., and Mihaescu, M.
\textit{Effects of tongue-base induced airway constriction on non-apneic palatal snoring}.
Coupled CFD--FSI, pressure/vorticity SPOD, and separation of hydrodynamic and acoustic modes.
\item \textbf{Li, P.}, Laudato, M., and Mihaescu, M.
\textit{Stiffness-dependent modal excitation and force--velocity coupling in fluid-induced soft-palate vibration}.
Modal coupling, nonlinear vibration, and sound-source analysis.
\end{enumerate}

\section{Additional Research Experience}
\textbf{Research Assistant -- Reduced Combustion Modelling} \hfill Sep 2020 -- Aug 2021\\
\textit{Institute for Technical Thermodynamics, Karlsruhe Institute of Technology, Germany}\\
Applied the reaction--diffusion manifold (REDIM) method to premixed flame--wall interactions;
validated reduced kinetic models against detailed kinetics.

\vspace{4pt}
\textbf{Research Project -- Low-Mach-Number Flow Modelling} \hfill Aug -- Sep 2018\\
\textit{Institute for Fluid Dynamics, Technical University of Darmstadt, Germany}\\
Used COMSOL FEM to investigate non-isothermal Couette flow and natural convection,
comparing low-Mach-number and Boussinesq formulations.

\vspace{4pt}
\textbf{Research Assistant -- Analytical and Numerical Solid Mechanics} \hfill Jul 2013 -- Jun 2015\\
\textit{Institute for Solid Mechanics, Dalian University of Technology, China}\\
Derived analytical plate bending and buckling solutions and compared them with FEM
using Mathematica, ABAQUS, and MATLAB.

\section{Scientific Communication}
Co-authored \textit{Effects of gravity force and tissue stiffness on a simplified
three-dimensional human soft palate}, ICSV30, Amsterdam (2024).
Conference participation: EFDC2, Dublin (2025); Swedish Mechanics Days, Gothenburg (2024);
Swedish Society of Biomechanics, Lund (2024).

\section{Languages}
Chinese -- Native \quad | \quad English -- IELTS 6.5/9.0 \quad | \quad
German -- TestDaF 17/20\\
\textbf{Swedish -- A2 (CEFR)}

\end{document}
